% ==================================================================
% id:       cond-021
% topic:    הסתברות מותנית ואי-תלות
% subtopic: מאורעות בלתי תלויים
% solution: full
% ==================================================================
שלושה ציידים יוצאים לצוד, וכל אחד מהם יורה פעם אחת. ההסתברות שצייד $A$ פוגע במטרה היא $\frac13$, וההסתברות שצייד $B$ פוגע בה היא $\frac14$. ההסתברות שלפחות אחד מהציידים פוגע במטרה היא $0.6$. הפגיעות של הציידים בלתי תלויות זו בזו.

מה ההסתברות שצייד $C$ פוגע במטרה?
\hint{עברו למשלים: „אף צייד לא פגע”. מה ההסתברות שלו, לפי אי־התלות?}
\begin{solbox}
נסמן ב־$A,B,C$ את המאורעות שכל צייד פוגע, וב־$p=\PP(C)$. המשלים של „לפחות אחד פגע” הוא „אף אחד לא פגע”, כלומר $A^C\cap B^C\cap C^C$. המשלימים של מאורעות בלתי תלויים הם בלתי תלויים, ולכן
\[
0.4=\PP(A^C\cap B^C\cap C^C)=\frac23\cdot\frac34\cdot(1-p)=\frac12(1-p).
\]
מכאן $1-p=0.8$, כלומר $p=\ans{0.2}$.
\end{solbox}
